\section*{Chapter \ref{ch:g12:phillips-curve} --- The Phillips Curve and Its Breakdown}
% figdata: figdata/grade-12/12-phillips-curve/out/ (phil6069, phil7083, phil1525, volcker, phillips_key).

\begin{solution}{exo:g12:phillips-curve:1}
1961: unemployment 6.7\,\%, inflation 1.1\,\%. 1969: 3.5\,\% and 5.5\,\%.
\end{solution}

\begin{solution}{exo:g12:phillips-curve:2}
1980 had both higher unemployment (7.1 against 4.5\,\%) and far higher inflation (13.5 against 1.6\,\%). On the 1960s line, unemployment of 7.1\,\% would come with inflation
near or below 1\,\%: 1980 lies far above it.
\end{solution}

\begin{solution}{exo:g12:phillips-curve:3}
$292.65/270.97 - 1 \approx 8.0\,\%$.
\end{solution}

\begin{solution}{exo:g12:phillips-curve:4}
At $u = 5$\,\%, the natural rate. At $u = 4$: $3 + 0.5 = 3.5\,\%$.
\end{solution}

\begin{solution}{exo:g12:phillips-curve:5}
$(2 + 3 + 1)/4 = 1.5$ point-years per point.
\end{solution}

\begin{solution}{exo:g12:phillips-curve:6}
Each year inflation rises by $0.5 \times 1 = 0.5$: 2.5, 3 and 3.5\,\%.
\end{solution}

\begin{solution}{exo:g12:phillips-curve:7}
Each year it falls by 0.5: 4.5, 4 and 3.5\,\%.
\end{solution}

\begin{solution}{exo:g12:phillips-curve:8}
With $u = u^n$, inflation stays at last year's, 3.5\,\%. Three years of unemployment one point lower were gained; the lasting cost is inflation 1.5 points higher, which only
a period of unemployment above the natural rate would bring down.
\end{solution}

\begin{solution}{exo:g12:phillips-curve:9}
Unemployment above 6\,\%: $1.1 + 1.6 + 3.7 + 3.6 = 10.0$ point-years; inflation fell from 13.5 to 3.2\,\%, by 10.3 points: a ratio of about 0.97.
\end{solution}

\begin{solution}{exo:g12:phillips-curve:10}
2019: unemployment 3.7\,\%, inflation 1.8\,\%. 2022: 3.6\,\% and 8.0\,\%. Almost the same unemployment came with very different inflation: the curve had shifted.
\end{solution}

\begin{solution}{exo:g12:phillips-curve:11}
With $\pi^e = 2$ every year, $\pi = 2 + 0.5 = 2.5\,\%$ each year: higher, but not accelerating.
\end{solution}

\begin{solution}{exo:g12:phillips-curve:12}
The claim: there is one fixed relation, so a permanently lower unemployment rate can be bought with permanently higher inflation. The evidence: in 1970--83 the United States had
both higher inflation and higher unemployment than in the 1960s (1975: 8.5\,\% and 9.1\,\%; 1980: 7.1\,\% and 13.5\,\%), so the relation moved up as inflation came to be expected.
\end{solution}

\begin{solution}{exo:g12:phillips-curve:13}
That year $\pi = 2 + 0 + 3 = 5\,\%$. The next, expectations are 5\,\% and $s = 0$: $\pi = 5\,\%$ again, and it stays there. To bring it back to 2\,\%, unemployment must be above 5\,\%
for a time: at 7\,\% it falls by one point a year, so three years.
\end{solution}

\begin{solution}{exo:g12:phillips-curve:14}
$\pi_t = \pi_0 + t\beta x$ grows without limit as $t$ grows. Holding $u$ below $u^n$ therefore means ever-rising inflation, which no central bank can accept for long: in the long
run unemployment returns to $u^n$ whatever the inflation rate, and the long-run curve is vertical.
\end{solution}

\begin{solution}{exo:g12:phillips-curve:15}
Expectations were \omterm{def:g12:phillips-curve:anchored}{anchored} near the 2\,\% target, so low unemployment did not feed into expected inflation; and the short-run slope $\beta$ may be small. Also, the natural rate
itself may have been lower than thought, so unemployment of 3.7\,\% was not far below it.
\end{solution}

\begin{solution}{pb:g12:phillips-curve:1}
\textbf{1.} Inflation: 1980, 13.5\,\%. Unemployment: 1982, 9.7\,\%.

\textbf{2.} $13.5 - 1.9 = 11.6$ points.

\textbf{3.} $13.5 - 3.2 = 10.3$ points.

\textbf{4.} 11.3\,\%.

\textbf{5.} No: 7.0\,\%.

\textbf{6.} $1.1 + 1.6 + 3.7 + 3.6 + 1.5 + 1.2 + 1.0 = 13.7$ point-years.

\textbf{7.} $13.7/11.6 \approx 1.18$.

\textbf{8.} Excesses $1.6 + 2.1 + 4.2 + 4.1 + 2.0 + 1.7 + 1.5 = 17.2$: ratio $17.2/11.6 \approx 1.48$.

\textbf{9.} Excesses $0.6 + 1.1 + 3.2 + 3.1 + 1.0 + 0.7 + 0.5 = 10.2$: ratio $10.2/11.6 \approx 0.88$.

\textbf{10.} Half a point on the natural rate moves the ratio by about 0.3: the ratio is only as precise as the natural rate, which is estimated.

\textbf{11.} $6.1 - 10.3 = -4.2 = -\beta \times 3.7$: $\beta \approx 1.14$.

\textbf{12.} $6.1 - 1.14 \times 3.6 \approx 2.0\,\%$, against 3.2\,\% actual.

\textbf{13.} $3.2 - 1.14 \times 1.5 \approx 1.5\,\%$, against 4.3\,\% actual: inflation rose while unemployment was above 6\,\%. Expectations did not simply follow last year's inflation,
the natural rate may have been higher than 6\,\%, or other shocks moved prices.

\textbf{14.} If people believe the announcement, expected inflation falls at once, and inflation can fall with less unemployment: a smaller ratio.

\textbf{15.} Quick: from 13.5 to 3.2\,\% in three years. Across 65 \omterm{def:g12:phillips-curve:disinflation}{disinflations}, quick ones had smaller (measured there as lost output) \omterm{def:g12:phillips-curve:disinflation}{sacrifice ratios}.

\textbf{16.} Mostly those who lost their jobs or could not find one, and their families.

\textbf{17.} Without a visible change of course, wage and price setters would have gone on expecting high inflation, and a larger rise in unemployment would have been needed.

\textbf{18.} Unemployment barely rose above its earlier level, so the point-years are close to zero; the fall in inflation came mostly from supply shocks fading, with expectations
already \omterm{def:g12:phillips-curve:anchored}{anchored}, not from slack.

\textbf{19.} The expectations-augmented curve: inflation fell when unemployment was above the natural rate, as it predicts, but its year-by-year forecasts with \omterm{def:g12:phillips-curve:adaptive}{adaptive
expectations} miss, as in 1984.

\textbf{20.} About 1.2 point-years of unemployment above 6\,\% per point of inflation removed (13.7 point-years for 11.6 points); between 0.9 and 1.5 for natural rates of
6.5\,\% and 5.5\,\%.
\end{solution}
